\section{Phase 3: Riemannian Diffusion Analysis of Structural and Geometric Preconditions}

Phase~3 investigates the structural and geometric conditions under which
compositional generative modelling becomes semantically meaningful or inevitably
fails. Unlike Phases~1 and~2, which evaluate the behaviour of compositional
operators directly, Phase~3 focuses on understanding \emph{why} these operators
succeed in some settings and break down in others. This phase directly supports
\textbf{Research Goal~3} and tests \textbf{Hypothesis~H3}, which proposes that
compositional correctness depends on latent-space alignment, coherent factor
interaction, and compatibility of the geometric structure of the underlying
densities.

To analyse these structural prerequisites, we employ \emph{Riemannian diffusion
models} (RDMs), which reinterpret diffusion as occurring on a curved data
manifold rather than in Euclidean pixel space. RDMs provide a principled way to
examine the geometry of the representations learned by the TB and Silicosis
models, enabling us to evaluate whether the manifolds underlying the two
distributions are compatible for composition. This geometric perspective offers
a diagnostic tool that is unavailable in standard diffusion frameworks.

\paragraph{Manifold Structure and Latent Alignment.}
Using RDMs, we examine whether the latent spaces of the TB and Silicosis models
share a compatible manifold structure. If the two diseases occupy regions with
aligned tangent spaces and similar local geometry, composition is expected to be
semantically meaningful. If their manifolds intersect at sharp angles, differ in
curvature, or exhibit incompatible local volume structure, composition is
expected to produce artifacts. This analysis provides a direct test of the first
component of Hypothesis~H3: that aligned or coherent latent structure is a
precondition for meaningful composition.

\paragraph{Compatibility of Semantic Directions.}
RDMs allow us to evaluate whether meaningful semantic transitions—such as
increasing cavitation severity or transitioning from nodular silicosis to
fibrotic TB—follow geodesics of the learned manifold. If semantic progression is
aligned with geodesic structure, compositional transitions have a natural
geometric interpretation. If semantic changes move off-manifold or violate
geodesic paths, composition will be structurally unstable regardless of the
operator used. This directly informs whether semantic factors are sufficiently
locally disentangled for compositional correctness.

\paragraph{Geometric Compatibility of Log-Densities.}
RDMs provide a framework for analysing whether diffusion models for TB and
Silicosis impose compatible geometric curvature on their log densities.
Compositional operators—particularly Product-of-Experts and density-based
methods like SuperDiff—implicitly assume that the geometry of the component
densities is consistent. RDMs allow us to visualise and characterise
incompatibilities: for example, when curvature induced by one pathology opposes
or overwhelms curvature induced by the other. Such conflict predicts failure
modes observed in previous phases. This analysis tests the final component of
Hypothesis~H3: that compositional correctness requires geometric compatibility
across component models.

\paragraph{Clinical Interpretation of Manifold Compatibility.}
We overlay the geometric findings with radiological interpretation to assess
whether regions of manifold compatibility correspond to clinically plausible
co-morbidity. When RDM analysis identifies smooth, aligned manifold regions,
compositions in those regions tend to reflect realistic hybrid pathology, such
as cavitations appearing alongside silicotic nodules. When RDMs reveal manifold
misalignment or curvature conflict, composed samples often contain unrealistic
hybrid textures, incorrect anatomical placement, or inconsistent opacity
patterns. This links geometric analysis directly to clinical semantics.

\paragraph{Contribution of This Phase.}
Phase~3 contributes a principled geometric and structural explanation for the
successes and failures of compositional generative modelling observed in the
earlier phases. By applying Riemannian diffusion analysis to TB and Silicosis
models, this phase identifies the latent and geometric prerequisites for valid
composition, clarifies the structural origins of failure modes, and offers a
unified geometric framework for understanding compositional operators across
score-based, PoE-based, and density-based methods. This completes the
methodological foundation for interpreting compositionality in complex domains.
